% V033 例题与变式 | 保留 examplebox 环境
\begin{examplebox}
\textbf{\textcolor{CExample}{例1：初中起点，四种辅助线分别计算}}
对照方法二至五的构造图。给定 $O=(0,0)$、$A=(3,4)$、$B=(1,2)$，求面积。
\begin{itemize}[leftmargin=*,itemsep=0.2em]
\item \textbf{过 $B$ 作竖直线：} $E_v=(1,4/3)$，$BE_v=2/3$；
$S=\frac12\cdot\frac23\cdot(1+2)=1$。
\item \textbf{过 $B$ 作水平线：} $E_h=(3/2,2)$，$BE_h=1/2$；
$S=\frac12\cdot\frac12\cdot(2+2)=1$。
\item \textbf{过 $A$ 作竖直线：} $F_v=(3,6)$，$AF_v=2$；
$S=\frac12\cdot2\cdot(3-2)=1$。
\item \textbf{过 $A$ 作水平线：} $F_h=(2,4)$，$AF_h=1$；
$S=\frac12\cdot1\cdot(4-2)=1$。
\item \textbf{最快核对：} $S=\frac12|3\cdot2-1\cdot4|=\boxed{1}$。
\end{itemize}
\textbf{识别信号：} 过 $B$ 的两辅助点在线段 $OA$ 内，用面积和；过 $A$ 的两辅助点在 $OB$ 延长线上，用面积差。

\medskip
\textbf{\textcolor{CExample}{例2：水平法遇到延长线——为什么必须相减？}}
$O=(0,0)$、$A=(4,3)$、$B=(2,-1)$。
\begin{itemize}[leftmargin=*,itemsep=0.2em]
\item \textbf{竖直法：} $E_v=(2,3/2)$，$BE_v=5/2$，
$S=\frac12\cdot\frac52\cdot(2+2)=5$。
\item \textbf{水平法：} $E_h=(-4/3,-1)$ 在 $OA$ 反向延长线上。
$BE_h=10/3$，$\triangle E_hAB$ 的高为 $4$，$\triangle E_hOB$ 的高为 $1$，
\[
S=\tfrac12\cdot\tfrac{10}{3}\cdot(4-1)=\boxed{5}.
\]
\item \textbf{快检：} $\tfrac12|4(-1)-2(3)|=5$。
\end{itemize}
\textbf{认识：} 辅助点在延长线上，若套“两个面积相加”会误判。

\medskip
\textbf{\textcolor{CExample}{例3：任意三点先平移，避免硬求高}}
$P=(1,2)$，$Q=(4,5)$，$R=(2,7)$。取 $P$ 为新原点，
$\overrightarrow{PQ}=(3,3)$，$\overrightarrow{PR}=(1,5)$，
\[
 S_{\triangle PQR}=\frac12|3\cdot5-3\cdot1|=\boxed{6}.
\]

\medskip
\end{examplebox}
\clearpage
\begin{examplebox}
\textbf{\textcolor{CExample}{例4：面积反求参数（绝对值两个分支）}}
$O=(0,0)$、$A=(t,3)$、$B=(2,4)$，且 $S_{\triangle OAB}=5$，求 $t$。
\begin{center}
\begin{tikzpicture}[x=1cm,y=1cm,font=\small,>=Stealth]
\path[use as bounding box] (-1,-1.6) rectangle (11,5.5);
\node[font=\small\bfseries,text=CStatement] at (5,5.1) {固定 OB，面积相同的 A 分处直线 OB 两侧};
\foreach \shift/\t/\sign/\sx in {0/4/+10/2.2,6.5/-1/-10/0.25}{
\begin{scope}[shift={(\shift+0.7,0)},scale=0.85]
\draw[->,gray] (-1.7,0)--(4.9,0) node[right] {$x$};
\draw[->,gray] (0,-0.4)--(0,4.8) node[above] {$y$};
\fill[CStatement!17] (0,0)--(\t,3)--(2,4)--cycle;
\draw[thick] (0,0)--(\t,3)--(2,4)--cycle;
\draw[very thick,CExample] (0,0)--(2,4);
\draw[dashed,CProof] (-1.4,3)--(4.5,3);
\node[below left] at (0,0) {$O$};
\node[above right] at (2,4) {$B(2,4)$};
\node[above] at (\t,3) {$A(\t,3)$};
\node at (\sx,2.5) {$S=5$};
\node at (2,5.3) {$4t-6=\sign$};
\end{scope}}
\node at (5,-1.15) {两图同尺度；$A$ 都在 $y=3$ 上，$t=4$ 或 $t=-1$。};
\end{tikzpicture}
\end{center}

\[
 \frac12|4t-6|=5\iff 4t-6=\pm10\iff
 \boxed{t=4\ \text{或}\ t=-1}.
\]

\medskip
\textbf{\textcolor{CExample}{例5：共线无需预检}}
$O=(0,0)$、$A=(2,3)$、$B=(4,6)$，
$S=\tfrac12|2\cdot6-4\cdot3|=0$。
\textbf{结论：} 面积公式本身就能检查三点共线。

\medskip
\textbf{\textcolor{CExample}{例6：由 V033 直接迁移到 C088}}
双曲线 $x^2/9-y^2/4=1$ 的渐近线为 $y=\pm\frac23x$。
选 $A=(6,4)$、$B=(3,-2)$，则
\[
 S_{\triangle OAB}=\frac12|6(-2)-3\cdot4|=12
 =\frac23|6\cdot3|.
\]
\textbf{迁移原则：} 把 $y_1=(b/a)x_1$、$y_2=-(b/a)x_2$ 直接代回 V033，得到 $S=\frac ba|x_1x_2|$。

\medskip
\textbf{\textcolor{CExample}{例7：变式（一般过原点两直线，抽象出斜率差）}}
$A$ 在 $y=m_1x$ 上，$B$ 在 $y=m_2x$ 上，$O=(0,0)$，
则 $A=(x_1,m_1x_1)$、$B=(x_2,m_2x_2)$。
\[
 \boxed{S_{\triangle OAB}=\frac12|m_2-m_1|\,|x_1x_2|.}
\]
\textbf{高价值迁移：} C088 是取 $m_1=b/a$、$m_2=-b/a$ 的特例。
\end{examplebox}
